\chapter*{คำนำ}
\addcontentsline{toc}{chapter}{คำนำ}

ตำราวิชาการเรื่อง \textbf{“เกษตรอัจฉริยะและปัญญาประดิษฐ์เพื่อการจัดการสวนทุเรียนและคุณภาพดิน” (Smart Agriculture and AIoT for Precision Durian Farming and Soil Analytics)} เล่มนี้ จัดทำขึ้นเพื่อเป็นเอกสารวิชาการหลักสำหรับการเรียนการสอน การวิจัย และการบริการวิชาการถ่ายทอดเทคโนโลยีสู่เกษตรกรชาวสวนไม้ผล โดยเฉพาะอย่างยิ่งเกษตรกรผู้ปลูกทุเรียนในภาคตะวันออก ซึ่งเป็นพืชเศรษฐกิจสำคัญอันดับหนึ่งของประเทศไทย

การเพาะปลูกทุเรียนในปัจจุบันกำลังเผชิญกับความท้าทายอย่างรุนแรงจากความแปรปรวนของสภาพภูมิอากาศ ปัญหาโรครากเน่าโคนเน่า การขาดแคลนแรงงาน และต้นทุนปัจจัยการผลิตที่เพิ่มสูงขึ้น การบริหารจัดการสวนแบบดั้งเดิมที่พึ่งพาเพียงประสบการณ์ไม่เพียงพอต่อการรักษาคุณภาพผลผลิตอีกต่อไป การบูรณาการองค์ความรู้ทางด้านฟิสิกส์เกษตร เครือข่ายเซนเซอร์อินเทอร์เน็ตของสรรพสิ่ง (Internet of Things; IoT), สถาปัตยกรรมสื่อสารระยะไกลกำลังต่ำ (LoRaWAN), และปัญญาประดิษฐ์เพื่อการประเมินสุขภาพดินจากภาพถ่ายดิจิทัล (SoilAI) จึงเป็นกุญแจสำคัญสู่การทำเกษตรแม่นยำสูง (High-Precision Agriculture)

เนื้อหาภายในเล่มแบ่งออกเป็น 8 บท ครอบคลุมตั้งแต่หลักการพื้นฐาน สถาปัตยกรรมระบบ ฟิสิกส์การเจริญเติบโตของผลทุเรียน การตรวจวัดสภาพแวดล้อม การควบคุมการให้น้ำอัจฉริยะ การวิเคราะห์คุณสมบัติดิน จนถึงกรณีศึกษาการนำไปใช้งานจริงในแปลงวิจัยและสวนเกษตรกรในจังหวัดจันทบุรี ผู้เขียนหวังเป็นอย่างยิ่งว่าตำราเล่มนี้จะเป็นประโยชน์อย่างยิ่งต่อนักศึกษา นักวิชาการ และเกษตรกรยุคใหม่ในการยกระดับภาคการเกษตรไทยสู่ความยั่งยืนในระดับสากล

\vspace{1.5cm}
\begin{flushright}
    \textbf{ผู้ช่วยศาสตราจารย์ ดร.ชีวะ ทัศนา}\\[0.2cm]
    สาขาวิชาฟิสิกส์ คณะวิทยาศาสตร์และเทคโนโลยี\\[0.1cm]
    มหาวิทยาลัยราชภัฏรำไพพรรณี\\[0.1cm]
    พุทธศักราช 2569
\end{flushright}
